\documentclass[18pt]{article}
\usepackage[utf8]{inputenc}
\usepackage[T1]{fontenc}
\usepackage{ragged2e}
\usepackage{caladea}
\usepackage{graphicx}
\usepackage{longtable}
\usepackage{wrapfig}
\usepackage{rotating}
\usepackage{epigraph}
\usepackage[normalem]{ulem}
\usepackage{hyperref}
\usepackage{amsmath}
\usepackage{amssymb}
\usepackage{capt-of}
\usepackage{hyperref}
\usepackage{fancyhdr}

\title{
 \includegraphics[scale=0.15, trim={10cm, 0, 10cm, 0}]{./assets/imagem.jpg}
  \par
   NOVENA DO LUTO}
\date{}

% Comando para fazer "Sumário" não aparecer no Sumário.
\renewcommand{\contentsname}{Sumário}
\begin{document}
\maketitle

\thispagestyle{empty} %zera a primeira página

\pagestyle{fancy}
\fancyhf{} % clear existing header/footer entries
\fancyfoot[LO, CE]{
  \includegraphics[scale=0.2]{./assets/cross.png} Ó Jesus, eu me abandono ao Senhor!
}
% Place Page X of Y on the right-hand
% side of the footer
\fancyfoot[R]{\thepage}

\newpage

\tableofcontents

\centering
\vfill
Visite-nos no Telegram: \url{https://t.me/CotidieNovena}
\newpage

\newpage


\begin{justify}

 \begin{center}
  \section{Apresentação}\label{sec:Apresentação} % (fold)
 \end{center}


 Durante os momentos de luto, quando o peso da perda parece insuportável, nossa fé pode se tornar uma luz guia na escuridão da dor. Esta novena é dedicada àqueles que navegam pelas etapas do luto após a morte de entes queridos. Ao longo desses nove dias de oração, vamos rezar para que você encontre consolo na presença de Deus e Sua graça enquanto avança pelas etapas do luto em direção à cura e à esperança.

 Na Igreja Católica Apostólica Romana, não há um santo patrono específico designado exclusivamente para o luto. No entanto, existem vários santos que são frequentemente invocados durante momentos de perda e dor. Esses santos são conhecidos por suas próprias lutas na vida e suas experiências de luto. Aqui está uma breve lista de alguns santos que são frequentemente considerados patronos para aqueles que estão de luto. Esta é apenas uma lista sugerida. Sinta-se à vontade para pedir a intercessão de qualquer santo para esta novena ou você pode simplesmente rezar a Deus.

\begin{itemize}
 \item \textbf{São José:} Dia da Festa: 19 de março. Patrono da Igreja universal, famílias, pais, trabalhadores e uma morte feliz.
 \item \textbf{Santa Rita de Cássia:} Dia da Festa: 22 de maio. Patrona de causas impossíveis, casamentos difíceis e vítimas de abuso.
 \item \textbf{São Judas Tadeu:} Dia da Festa: 28 de outubro. Patrono de casos desesperados, situações desesperadoras e causas perdidas.
 \item \textbf{Nossa Senhora das Dores:} Dia da Festa: 15 de setembro. Patrona dos que sofrem, mães e os que estão de luto.
 \item \textbf{Santa Isabel da Hungria:} Dia da Festa: 17 de novembro. Patrona de padeiros, hospitais, pessoas sem-teto e viúvas.
 \item \textbf{Santa Elizabeth Ann Seton:} Dia da Festa: 4 de janeiro. Patrona de escolas católicas, viúvas e perda de pais.
 \item \textbf{Santa Gianna Molla:} Dia da Festa: 28 de abril. Patrona de médicos, mães, crianças não nascidas e o movimento pró-vida.
\end{itemize}

 Você pode rezar a Novena do Luto por qualquer intenção, mas especialmente se você ou alguém que você conhece está sofrendo a dor da morte de um membro da família, amigo ou ente querido.

\end{justify}

%%%%%%%%%%%%%%%%%%%%%%%%%%%%%%%%%%%%% Orações  %%%%%%%%%%%%%%%%%%%%%%%%%%%%%%%%%%%%%%%%%%%

\newpage
\section{Orações}\label{sec:Orações} % (fold)
\subsection{Oração Inicial}\label{sec:Oração_Inicial} % (fold)

Pai Celestial, reconheço a fase de negação que muitas vezes acompanha o luto e a perda. Nestes momentos em que não consigo acreditar, quando a realidade parece distante, volto-me para Ti em busca de conforto. Ajuda-me a enfrentar minhas emoções com coragem, enquanto me apoio no Teu amor para me guiar através desta fase inicial de tristeza e negação. Senhor, humildemente Te peço que ouças minhas intenções.

(Mencione suas intenções aqui.)

Por favor, seja meu refúgio na confusão, meu consolo na incerteza e minha força enquanto começo a entender a dolorosa ausência dos meus entes queridos. Amém.

Salve, Rainha, Mãe de Misericórdia, nossa vida, nossa doçura e nossa esperança.

A Ti clamamos, pobres filhos de Eva.

A Ti suspiramos, gemendo e chorando neste vale de lágrimas.

Vendei, pois, ó Advocada nossa, os olhos de misericórdia para nós, e depois deste nosso exílio, mostrai-nos o bendito fruto do vosso ventre, Jesus.

Ó clemente, ó amorosa, ó doce Virgem Maria. Amém.

\subsection{Novena do Luto - Dia 2}
Querido Senhor, em meio à tempestade da raiva que muitas vezes surge na experiência do luto, humildemente peço a Tua graça para encontrar paz em meu coração, em minha mente e em meu corpo. Enquanto luto com as emoções intensas que surgem da sensação de injustiça e perda que sinto, ajuda-me a canalizar minha raiva de forma construtiva. Senhor, humildemente Te peço que ouças minhas intenções.

(Mencione suas intenções aqui.)

Que minhas frustrações sejam transformadas em uma busca por compreensão e aceitação. Por favor, enche meu coração com a paz que somente a Tua presença pode proporcionar. Luto com o desejo em meu coração de me reunir novamente com meu ente querido. Guia-me para encontrar paz em meio à turbulência da minha raiva e emoções. Amém.

Salve, Rainha, Mãe de Misericórdia, nossa vida, nossa doçura e nossa esperança.

A Ti clamamos, pobres filhos de Eva.

A Ti suspiramos, gemendo e chorando neste vale de lágrimas.

Vendei, pois, ó Advocada nossa, os olhos de misericórdia para nós, e depois deste nosso exílio, mostrai-nos o bendito fruto do vosso ventre, Jesus.

Ó clemente, ó amorosa, ó doce Virgem Maria. Amém.

\subsection{Novena do Luto - Dia 3}
Deus amoroso, durante a fase de barganha, quando me encontro buscando uma maneira de desfazer o irreversível, por favor, abrace-me! Atrai-me para o Teu abraço curador! Senhor, humildemente Te peço que ouças minhas intenções.

(Mencione suas intenções aqui.)

Trago meus pensamentos sobre o que poderia ter sido e o que eu gostaria de poder mudar ou ter feito de forma diferente diante de Ti, ó Senhor. Por favor, levante o peso da culpa e a dor do "se ao menos". Ensina-me a entregar meus desejos por resultados alternativos e a abraçar a beleza das vidas que foram vividas, não importa quão longas ou curtas tenham sido. Assegura-me que as almas dos meus entes queridos estão agora em Tuas mãos cuidadosas. Amém.

Salve, Rainha, Mãe de Misericórdia, nossa vida, nossa doçura e nossa esperança.

A Ti clamamos, pobres filhos de Eva.

A Ti suspiramos, gemendo e chorando neste vale de lágrimas.

Vendei, pois, ó Advocada nossa, os olhos de misericórdia para nós, e depois deste nosso exílio, mostrai-nos o bendito fruto do vosso ventre, Jesus.

Ó clemente, ó amorosa, ó doce Virgem Maria. Amém.

\subsection{Novena do Luto - Dia 4}
Senhor da Esperança, enquanto atravesso as profundezas da depressão que muitas vezes acompanham o l uto, peço-Te que me enchas de esperança que transcende e supera o desespero. Quando o peso da tristeza parece insuportável, quando sinto que não consigo respirar, ficar de pé ou até mesmo falar, levanta-me do abismo de dor e lágrimas e guia-me em direção à luz da Tua promessa. Senhor, humildemente Te peço que ouças minhas intenções.

(Mencione suas intenções aqui.)

Que eu encontre consolo e esperança ao saber que Tu recolhes cada uma das minhas lágrimas e que há um amanhecer de alegria além das noites de choro. Que meu coração seja renovado com a expectativa do futuro onde a dor não existirá mais e a vida será cheia de esperança. Amém.

Salve, Rainha, Mãe de Misericórdia, nossa vida, nossa doçura e nossa esperança.

A Ti clamamos, pobres filhos de Eva.

A Ti suspiramos, gemendo e chorando neste vale de lágrimas.

Vendei, pois, ó Advocada nossa, os olhos de misericórdia para nós, e depois deste nosso exílio, mostrai-nos o bendito fruto do vosso ventre, Jesus.

Ó clemente, ó amorosa, ó doce Virgem Maria. Amém.

\subsection{Novena do Luto - Dia 5}
Pai Celestial, à medida que a fase da aceitação gradualmente surge em meu coração, busco a Tua orientação para abraçar as memórias e experiências que definem as vidas dos meus entes queridos. Por favor, dá-me a graça de ser grato pelo tempo que compartilhamos, não importa quão longo ou curto tenha sido, e pelo amor que vivenciamos. Senhor, humildemente Te peço que ouças minhas intenções.

(Mencione suas intenções aqui.)

Nesta fase de luto e dor, ajuda-me a encontrar paz na aceitação da ausência física deles. Por favor, ajuda-me a valorizar sua presença em meu coração para sempre. Ensina-me que aceitar não é esquecer, mas um passo em direção à cura da dor da perda. Amém.

Salve, Rainha, Mãe de Misericórdia, nossa vida, nossa doçura e nossa esperança.

A Ti clamamos, pobres filhos de Eva.

A Ti suspiramos, gemendo e chorando neste vale de lágrimas.

Vendei, pois, ó Advocada nossa, os olhos de misericórdia para nós, e depois deste nosso exílio, mostrai-nos o bendito fruto do vosso ventre, Jesus.

Ó clemente, ó amorosa, ó doce Virgem Maria. Amém.

\subsection{Novena do Luto - Dia 6}
Querido Senhor, enquanto navego pela fase de encontrar significado em meu luto, ajuda-me a descobrir as lições e o propósito ocultos dentro da minha dor. Senhor, humildemente Te peço que ouças minhas intenções.

(Mencione suas intenções aqui.)

No meio do sofrimento, por favor, transforma minha tristeza em um catalisador para o crescimento pessoal – crescimento na fé, crescimento na empatia e crescimento na compaixão pelos outros. Que minha jornada através do luto, que dói tanto, não seja em vão, mas um caminho de transformação, convertendo minha dor em uma força para mudanças positivas em minha vida e uma fé mais profunda e dependência de Ti, ó Senhor, enquanto me sustentas nesta dolorosa experiência de luto. Amém.

Salve, Rainha, Mãe de Misericórdia, nossa vida, nossa doçura e nossa esperança.

A Ti clamamos, pobres filhos de Eva.

A Ti suspiramos, gemendo e chorando neste vale de lágrimas.

Vendei, pois, ó Advocada nossa, os olhos de misericórdia para nós, e depois deste nosso exílio, mostrai-nos o bendito fruto do vosso ventre, Jesus.

Ó clemente, ó amorosa, ó doce Virgem Maria. Amém.

\subsection{Novena do Luto - Dia 7}
Deus amoroso, na fase de luto de alcançar outros que compartilham minha dor, dá-me a força para me conectar mesmo que eu sinta que não sou forte o suficiente. Por favor, proporciona-me a capacidade de ser uma fonte genuína de conforto e apoio. Senhor, humildemente Te peço que ou ```tex
ças minhas intenções.

(Mencione suas intenções aqui.)

Por favor, usa minhas experiências de luto para aprofundar minha empatia. Faz de mim um instrumento da Tua compaixão curadora, para que outros possam ver o Teu rosto curador. À medida que estendo minha mão àqueles que precisam, que minha jornada se torne uma fonte de esperança, lembrando a todos que não estão sozinhos em sua dor. Amém.

Salve, Rainha, Mãe de Misericórdia, nossa vida, nossa doçura e nossa esperança.

A Ti clamamos, pobres filhos de Eva.

A Ti suspiramos, gemendo e chorando neste vale de lágrimas.

Vendei, pois, ó Advocada nossa, os olhos de misericórdia para nós, e depois deste nosso exílio, mostrai-nos o bendito fruto do vosso ventre, Jesus.

Ó clemente, ó amorosa, ó doce Virgem Maria. Amém.

\subsection{Novena do Luto - Dia 8}
Senhor de Todo o Tempo, durante a fase de entrega, ajuda-me a renunciar ou deixar ir minha necessidade de controle para que eu possa encontrar consolo em confiar no Teu tempo. Senhor, humildemente Te peço que ouças minhas intenções.

(Mencione suas intenções aqui.)

Por favor, fortalece minha fé enquanto coloco minha dor, tristeza, perguntas e dúvidas em Tuas mãos. Ensina-me que entregar-me à Tua santa vontade não é um sinal de fraqueza, mas um ato de humildade. Ajuda-me a saber que os Teus caminhos estão além da minha compreensão. À medida que me entrego, que eu experimente a paz que vem de confiar completamente em Ti. Amém.

Salve, Rainha, Mãe de Misericórdia, nossa vida, nossa doçura e nossa esperança.

A Ti clamamos, pobres filhos de Eva.

A Ti suspiramos, gemendo e chorando neste vale de lágrimas.

Vendei, pois, ó Advocada nossa, os olhos de misericórdia para nós, e depois deste nosso exílio, mostrai-nos o bendito fruto do vosso ventre, Jesus.

Ó clemente, ó amorosa, ó doce Virgem Maria. Amém.

\subsection{Novena do Luto - Dia 9}
Pai Celestial, ao concluir esta novena pelo luto, busco renovação em Ti e em minha fé na Tua esperança eterna. Que as etapas do luto se tornem degraus em direção à cura e ao crescimento. Senhor, humildemente Te peço que ouças minhas intenções.

(Mencione suas intenções aqui.)

Enquanto honro meus amados falecidos, por favor, guia-me em direção a um futuro enriquecido por suas memórias e sustentado pela Tua graça. Transforma minha dor e sofrimento em um testemunho de amor duradouro por eles. No meio da perda, que o Teu amor permaneça meu âncora constante, mantendo-me enraizado na minha fé em Ti, ó Senhor. Amém.

Salve, Rainha, Mãe de Misericórdia, nossa vida, nossa doçura e nossa esperança.

A Ti clamamos, pobres filhos de Eva.

A Ti suspiramos, gemendo e chorando neste vale de lágrimas.

Vendei, pois, ó Advocada nossa, os olhos de misericórdia para nós, e depois deste nosso exílio, mostrai-nos o bendito fruto do vosso ventre, Jesus.

Ó clemente, ó amorosa, ó doce Virgem Maria. Amém.

\subsection*{Créditos:}
\href{https://www.praymorenovenas.com/bereavement-novena}{Pray More Novenas}

\end{document}
